\begin{tikzpicture}[
  >={Latex[length=3mm,width=2.2mm]},
  every node/.style={font=\small},
  uvec/.style={->,line width=1pt,black},
  pvec/.style={->,line width=1.3pt,vpurple},
  posc/.style={circle,draw=black!60,fill=posred,inner sep=0pt,minimum size=4.5mm},
  negc/.style={circle,draw=black!60,fill=negblue,inner sep=0pt,minimum size=4.5mm}
]

\coordinate (O) at (0,0);
\coordinate (P) at (6,4.2);
\coordinate (qp) at (0,1.0);    % +q
\coordinate (qm) at (0,-1.0);   % -q

% Asse verticale
\draw[thin,black] (0,-1.8) -- (0,4.5);

% Segmento a
\draw[thin,black] (-1.4,1.0) -- (-1.4,-1.0);
\draw[thin,black] (-1.5,1.0) -- (-1.3,1.0);
\draw[thin,black] (-1.5,-1.0) -- (-1.3,-1.0);
\node[left] at (-1.4,0) {$a$};

% Cariche
\node[posc] at (qp) {\tiny$+$}; \node[left] at (-0.3,1.0) {$+q$};
\node[negc] at (qm) {\tiny$-$}; \node[left] at (-0.3,-1.0) {$-q$};
\node[left] at (0,0) {$O$};
\fill (O) circle (1pt);

% Dipolo p
\draw[pvec] (0,-1.0) -- (0,1.5) node[above left] {$\mathbf{p}$};

% Triangolo da +q a P (r_1), da -q a P (r_2), da O a P (r)
\draw[thin] (qp) -- node[midway,above] {$r_1$} (P);
\draw[thin] (qm) -- node[pos=0.5,below right] {$r_2$} (P);
\draw[thin] (O)  -- node[midway,above] {$r$} (P);

% u_r
\draw[uvec] (O) -- ++(0.9,0.65) node[right] {$\mathbf{u}_r$};

% Angolo theta
\draw[thin] (O) ++(0.3,0) arc (0:60:0.3);
\node at (0.4,0.35) {$\theta$};

% P
\fill (P) circle (1.2pt);
\node[right] at (P) {$P$};

% r_2 - r_1 (segmento proiettato sulla retta)
\draw[thin] (qm) -- (-0.01,-0.35);
\fill[ytangent,draw=black!40,opacity=0.6] (qm) -- (qp) -- (0,-0.3) -- cycle;
\node at (1.1,-0.3) {$r_2 - r_1$};

\end{tikzpicture}
